\documentclass{article}
\usepackage{graphicx}
\usepackage{amsmath}   % para ecuaciones
\usepackage{amssymb}   % para símbolos como \mathbb{R}
\usepackage{amsfonts}  % para fuentes matemáticas
\usepackage[margin=0.7in]{geometry}  % márgenes reducidos

\title{PS0 - Solutions}
\author{Facundo Berthet}
\date{April 2026}

\begin{document}
\maketitle

\section{Gradients and Hessians}

\subsection*{(a)}
Let $f(x) = \frac{1}{2}x^TAx + b^Tx$, where $A$ is symmetric and $b \in \mathbb{R}^n$. 
We compute the gradient by splitting into two terms:

\textbf{Term 1:} Using $\nabla_x x^TAx = 2Ax$ (since $A$ is symmetric):
$$\nabla_x\left(\frac{1}{2}x^TAx\right) = \frac{1}{2} \cdot 2Ax = Ax$$

\textbf{Term 2:} Using $\nabla_x b^Tx = b$:
$$\nabla_x\left(b^Tx\right) = b$$

\textbf{Therefore:}
$$\nabla f(x) = Ax + b$$

\subsection*{(b)}
Let $f(x) = g(h(x))$, where $g : \mathbb{R} \to \mathbb{R}$ and 
$h : \mathbb{R}^n \to \mathbb{R}$ are differentiable. We apply the chain rule.

For each partial derivative with respect to $x_i$:
$$\frac{\partial f}{\partial x_i} = g'(h(x)) \cdot \frac{\partial h(x)}{\partial x_i}$$

Stacking all partial derivatives into a vector:
$$\boxed{\nabla f(x) = g'(h(x)) \cdot \nabla h(x)}$$

where $g'(h(x)) \in \mathbb{R}$ is a scalar and $\nabla h(x) \in \mathbb{R}^n$ is a vector.

\subsection*{(c)}
Let $f(x) = \frac{1}{2}x^TAx + b^Tx$. From part (a), we know that 
$\nabla f(x) = Ax + b$. Taking the gradient again:

$$\nabla^2 f(x) = A$$

since $\nabla^2_x(Ax) = A$ and $\nabla^2_x(b) = 0$.

\subsection*{(d)}
Let $f(x) = g(a^Tx)$, where $g : \mathbb{R} \to \mathbb{R}$ and $a \in \mathbb{R}^n$.

\textbf{Gradient:} This is the chain rule from (b) with $h(x) = a^Tx$:
$$\nabla f(x) = g'(a^Tx) \cdot a$$

\textbf{Hessian:} We differentiate the gradient with respect to $x$. 
Since $a$ is constant, we apply the chain rule to $g'(a^Tx)$. 
The $(i,j)$ entry of the Hessian is:
$$\frac{\partial^2 f}{\partial x_i \partial x_j} = g''(a^Tx) \cdot a_i \cdot a_j$$

Recognizing that $a_i a_j$ is the $(i,j)$ entry of the outer product $aa^T$:
$$\nabla^2 f(x) = g''(a^Tx)(aa^T)$$


\section{Positive Definite Matrices}

\subsection*{(a)}
We want to show that $A = zz^T$ is PSD, i.e., $x^TAx \geq 0$ 
for all $x \in \mathbb{R}^n$. Substituting $A = zz^T$:

$$x^TAx = x^T(zz^T)x = (x^Tz)(z^Tx) = (x^Tz)^2 \geq 0$$

where the last inequality holds since any real number squared is 
non-negative. $\blacksquare$

\subsection*{(b)}
Let $A = zz^T$ with $z \neq 0$. 

\textbf{Nullspace:} We need $Ax = zz^Tx = z(z^Tx) = 0$. Since 
$z \neq 0$, this requires $z^Tx = 0$. Therefore:
$$\mathcal{N}(A) = \{x \in \mathbb{R}^n : z^Tx = 0\}$$
which is the set of all vectors orthogonal to $z$, a subspace 
of dimension $n-1$.

\textbf{Rank:} By the rank-nullity theorem:
$$\text{rank}(A) + \dim(\mathcal{N}(A)) = n 
\implies \text{rank}(A) = n - (n-1) = 1$$

\subsection*{(c)}
We claim $BAB^T$ is PSD. For any $x \in \mathbb{R}^m$, 
let $y = B^Tx \in \mathbb{R}^n$. Then:

$$x^T(BAB^T)x = (B^Tx)^T A (B^Tx) = y^TAy \geq 0$$

where the last inequality holds since $A$ is PSD by hypothesis. 
Since this holds for all $x \in \mathbb{R}^m$, $BAB^T$ is PSD. 
$\blacksquare$

\section{Eigenvectors, Eigenvalues, and the Spectral Theorem}

\subsection*{(a)}
Given $A = T\Lambda T^{-1}$, we multiply both sides on the right by $T$:
$$AT = T\Lambda$$

Looking at the $i$-th column of each side. The $i$-th column of 
the left side is $At^{(i)}$. Since multiplying by a diagonal matrix 
on the right scales each column, the $i$-th column of $T\Lambda$ is 
$\lambda_i t^{(i)}$. Therefore:
$$At^{(i)} = \lambda_i t^{(i)} \quad \blacksquare$$

\subsection*{(b)}
Given $A = U\Lambda U^T$, we multiply on the right by $u^{(i)}$:
$$Au^{(i)} = U\Lambda U^T u^{(i)}$$

Since the columns of $U$ are orthonormal, $U^Tu^{(i)} = e_i$, 
the $i$-th canonical vector:
$$(U^Tu^{(i)})_j = (u^{(j)})^T u^{(i)} = \begin{cases} 1 & j = i \\ 
0 & j \neq i \end{cases}$$

Applying $\Lambda$ and then $U$:
$$Au^{(i)} = U\Lambda e_i = U(\lambda_i e_i) = \lambda_i Ue_i = 
\lambda_i u^{(i)} \quad \blacksquare$$

\subsection*{(c)}
Since $A$ is PSD, by definition $x^TAx \geq 0$ for all 
$x \in \mathbb{R}^n$. Choosing $x = u^{(i)}$:
$$u^{(i)T}Au^{(i)} \geq 0$$

From part (b), $Au^{(i)} = \lambda_i u^{(i)}$, therefore:
$$u^{(i)T}(\lambda_i u^{(i)}) = \lambda_i \underbrace{u^{(i)T}
u^{(i)}}_{=1} = \lambda_i \geq 0 \quad \blacksquare$$

\end{document}